%% 2i2d-pv-associations — trois branchements de quatre cellules photovoltaïques (18 V ; 5,56 A).
%% Le branchement fixe le couple (U, I) mais pas leur produit : la puissance reste la somme des
%% puissances des cellules. Rail « + » en solideA, rail « - » en noir.
\documentclass[tikz,border=4pt]{standalone}
\usepackage{liaisons}
\begin{document}
\begin{tikzpicture}[x=1cm,y=1cm,
  film/.style={line width=0.8pt, line cap=round},
  filp/.style={line width=0.8pt, line cap=round, draw=solideA},
  titre/.style={font=\scriptsize\bfseries, anchor=south},
  val/.style={font=\scriptsize, anchor=north, inner sep=1pt},
  puis/.style={font=\scriptsize\bfseries, text=solideE, anchor=north, inner sep=2.5pt,
               draw=solideE, line width=0.7pt, rounded corners=2pt}]

%% ---- une cellule photovoltaïque, centrée en (#1,#2) ----------------------
\newcommand\cellule[2]{%
  \begin{scope}
    \path[clip] (#1-0.36,#2-0.31) rectangle (#1+0.36,#2+0.31);
    \fill[solideB!22] (#1-0.36,#2-0.31) rectangle (#1+0.36,#2+0.31);
    \foreach \k in {-3,-2,...,3} { \draw[solideB!55, line width=0.35pt]
      (#1+0.09*\k,#2-0.31) -- (#1+0.09*\k,#2+0.31); }
  \end{scope}
  \draw[solideB, line width=0.7pt] (#1-0.36,#2-0.31) rectangle (#1+0.36,#2+0.31);
  \node[font=\tiny, fill=white, inner sep=0.4pt] at (#1,#2+0.15) {18 V};
  \node[font=\tiny, fill=white, inner sep=0.4pt] at (#1,#2-0.15) {5,56 A};}
\newcommand\bornemoins[2]{\fill (#1,#2) circle (1.1pt);
  \node[font=\scriptsize, anchor=south] at (#1,#2+0.04) {$-$};}
\newcommand\borneplus[2]{\fill[solideA] (#1,#2) circle (1.1pt);
  \node[font=\scriptsize, text=solideA, anchor=south] at (#1,#2+0.02) {$+$};}

%% ================= 1. PARALLÈLE ==========================================
\node[titre] at (0.55,3.36) {Parallèle};
\foreach \y in {2.77,2.03,1.29,0.55} {
  \cellule{0.55}{\y}
  \draw[film] (0,\y) -- (0.19,\y);
  \draw[filp] (0.91,\y) -- (1.1,\y);}
\draw[film] (0,0.55) -- (0,3.00);
\draw[filp] (1.1,0.55) -- (1.1,3.00);
\bornemoins{0}{3.00}  \borneplus{1.1}{3.00}

%% ================= 2. SÉRIE-PARALLÈLE ====================================
\node[titre] at (2.55,3.36) {Série-parallèle};
\foreach \y in {2.05,1.25} {
  \cellule{2.10}{\y}  \cellule{3.00}{\y}
  \draw[film] (1.55,\y) -- (1.74,\y);
  \draw[line width=0.8pt] (2.46,\y) -- (2.64,\y);
  \draw[filp] (3.36,\y) -- (3.55,\y);}
\draw[film] (1.55,1.25) -- (1.55,3.00);
\draw[filp] (3.55,1.25) -- (3.55,3.00);
\bornemoins{1.55}{3.00}  \borneplus{3.55}{3.00}

%% ================= 3. SÉRIE ==============================================
\node[titre] at (6.0,3.36) {Série};
\foreach \x in {4.66,5.56,6.46,7.36} { \cellule{\x}{1.66} }
\draw[film] (4.1,1.66) -- (4.30,1.66);
\draw[line width=0.8pt] (5.02,1.66) -- (5.20,1.66)  (5.92,1.66) -- (6.10,1.66)
                        (6.82,1.66) -- (7.00,1.66);
\draw[filp] (7.72,1.66) -- (7.9,1.66);
\draw[film] (4.1,1.66) -- (4.1,3.00);
\draw[filp] (7.9,1.66) -- (7.9,3.00);
\bornemoins{4.1}{3.00}  \borneplus{7.9}{3.00}

%% ================= tableau de synthèse ===================================
\rappel{\foreach \x/\u/\i in {0.55/18/{22{,}24}, 2.55/36/{11{,}12}, 6.0/72/{5{,}56}} {
  \node[val] at (\x,0.02) {$U = \u$ V};
  \node[val] at (\x,-0.33) {$I = \i$ A};
  \node[puis] at (\x,-0.72) {$P \approx 400$ W};}

\draw[solideE, line width=0.7pt] (0.55,-1.26) -- (6.0,-1.26);
\foreach \x in {0.55,2.55,6.0} { \draw[solideE, line width=0.7pt] (\x,-1.26) -- (\x,-1.18); }
\draw[solideE, line width=0.7pt] (3.28,-1.26) -- (3.28,-1.34);
\node[font=\scriptsize, text=solideE, anchor=north] at (3.28,-1.34)
  {$P = U \times I$ : même puissance dans les trois cas};
\node[font=\scriptsize, anchor=north, text=black!65]
  at (3.28,-1.68)
  {valeur exacte 400,3 W, arrondie ici à 400 W};}
\end{tikzpicture}
\end{document}
